\subsection{Structural definition and function of the gravitational constant}
\label{sec:ii-definicion-gravedad}

\begin{definition}[Gravitational section contragredient to action]
For a positive section \(\hbar_a\), a velocity
\(c_{\rm int}>0\), and the chronological radius \(R_{108}\) of
section~\ref{sec:gravity}, the circular gravitational section is
\begin{equation}
 G_a=\frac{c_{\rm int}^{3}R_{108}^{2}}{\hbar_a},
 \qquad R_{108}=\frac{54}{\pi}\ell_0,
 \qquad \ell_0=c_{\rm int}t_0.
 \label{eq:ii-definicion-gravedad-seccion}
\end{equation}
It is the action-reciprocal section realizing coincidence
between the cycle's reduced gravitational radius and its
chronological radius, in the convention \(r_g=Gm/c_{\rm int}^2\).
\end{definition}

The radius-coincidence proposition of
section~\ref{sec:gravity} proves its uniqueness in the collection
\(G_a(\vartheta)=\vartheta c_{\rm int}^{5}t_0^2/\hbar_a\).
The circular condition fixes \(\vartheta=(54/\pi)^2\) and leads
to \eqref{eq:ii-definicion-gravedad-seccion}. That condition and the
positive dimensional sections are part of the realization.
Section~\ref{md:rigidez} constructs its length and brings together
the uniqueness proof for the general positive character, with
$R_{108}=L_\alpha$. The circular selector and longitudinal readout
operator therefore preserve the same normalization.
Keeping them common to both orientations yields
\begin{equation}
 \frac{G_{\mathrm{pre}}}{G_{\mathrm{ret}}}
 =R_{\mathrm{act}}^{-1},
 \qquad
 \frac{G_a\hbar_a}{c_{\rm int}^{3}}
 =R_{108}^{2}=\ell_P^2.
 \label{eq:ii-definicion-gravedad-area}
\end{equation}
The proof is direct substitution of
\eqref{eq:ii-definicion-gravedad-seccion}; the action change
is compensated by the reciprocal change in \(G_a\).
